% !PRINTED 156
\weberpage{156}

If in (15) we set $v=K$ and define a quantity $E'$ by
the equation:
\weq{17}{E(K+iK') = E + i(K'-E'),}
then one obtains the formula known under the name of \gesperrt{Legendre's relation}
\weq{18}{EK' + KE' - KK' = \frac{\pi}{2}\,.}

The significance of the quantity $E'$ introduced here is recognised
if one derives \gesperrt{Legendre's relation} by a second route,
making use of one of the linear transformations.

It follows, if one sets in \S\,31 (11):
\[
u = \frac{v}{2K}, \qquad \omega = \frac{iK'}{K}, \qquad \frac{u}{\omega} = \frac{v}{2iK'},
\]
by \S\,43 (1), (4):
\weq{19}{\sqrt{K}\, e^{-\frac{\pi v^2}{4KK'}} \,\Theta(iv,\varkappa') = \sqrt{K'}\, H(v+K),}
and from this, by logarithmic differentiation,
\weq{20}{iZ(iv,\varkappa') = \frac{\pi v}{2KK'} + \frac{H'(v+K)}{H(v+K)}}
or, by \S\,43 (6),
\weq{21}{iZ(iv,\varkappa') = \frac{\pi v}{2KK'} + Z(v) + \frac{d\log \mathrm{cn}\,v}{dv}\,.}

But by (5) we have
\[
i\,\frac{dZ(iv,\varkappa')}{dv} = -\,\mathrm{dn}^2(iv,\varkappa') + \frac{E'}{K'},
\]
if $E'$ now has the meaning:
% [sic] German source prints the lower limit v; for E' to be the complete integral obtained from E by interchanging \varkappa and \varkappa', the lower limit is mathematically expected to be 0.
\weq{22}{E' = \int\limits_v^{K'} \mathrm{dn}^2(v,\varkappa')\, dv,}
so that it arises from $E$ by the interchange of $\varkappa$ with $\varkappa'$. But if one
sets $v=0$ in (21), after having first differentiated, the relation (18)
again results, from which it follows that $E'$ is in both
cases the same quantity.

\begin{center}
\textbf{\S\,48. \gesperrt{The elliptic transcendents of the third kind.}}
\end{center}

If we take the formula (10) of the preceding section:
\[
\frac{\Theta'(u+v)}{\Theta(u+v)} - \frac{\Theta'(u-v)}{\Theta(u-v)} = 2Z(v) - \frac{2\varkappa^2\,\mathrm{sn}^2 u\,\mathrm{sn}\,v\,\mathrm{cn}\,v\,\mathrm{dn}\,v}{1 - \varkappa^2\,\mathrm{sn}^2 u\,\mathrm{sn}^2 v}
\]
%CONT>>
